\pagebreak

\chapter{Inbuilt Modules}

    \index{module!inbuilt}
    The following called modules are built into the compiler. User
    called modules or procedures must not use identifiers that conflict
    with these.
    
    The following abbreviations are used to denote modes -

    \begin{tabular}{| l | l |}
    \hline
    I  & immediate      \\ \hline
    IN & input          \\ \hline
    V  & value          \\ \hline
    SV & selectvalue    \\ \hline
    S  & static         \\ \hline
    Q  & queue          \\ \hline
    M  & memory         \\ \hline
    PR & priority       \\ \hline
    C  & clock          \\ \hline
    \end{tabular}
    
    The following abbreviations are used to denote types -

    \begin{tabular}{| l | l |}
    \hline
    bits & bits                                                 \\ \hline
    uint & unsigned integer                                     \\ \hline
    int & integer or unsigned integer                           \\ \hline
    fixed & fixed point                                         \\ \hline
    float & floating point                                      \\ \hline
    log & logical (boolean)                                     \\ \hline
    str & string                                                \\ \hline
    type & type                                                 \\ \hline
    any & any type                                              \\ \hline
    prim & a primitive type                                     \\ \hline
    comp & a compound type (arrays, structs or any combination) \\ \hline
    [] & array of type                                          \\ \hline
    \end{tabular}


    Modules with 'yes' to 'Parameter=argument allowed' can have arguments
    matched to parameters by position in the argument list or by
    parameter=argument (keymatch) or a mixture as long as positional arguments
    do not follow keymatch arguments. Otherwise these keymatch arguments are
    not allowed. This may be because - \\
    \blist
    \item   the module has a single input or output argument in which case
    no positional or keymatch parameter association is needed
    \item   the module has an indefinite number of arguments and hence there
    are no 'named' parameters with which to associate arguments
    \blend
    The keymatching may be restricted to just input or just output.


    \section{del - variable length delay line}\label{sec:imdel}
        
        {\bf Call:} \\
        del(out \verb'<-' in, len, init, step, varlen, reset);
        
        {\bf Output parameters:} \\
        \begin{tabular}{| l | l | l | l |}
        \hline
        param  & mode & type  & \\
        \hline \hline
        out  & V & any & output \\ \hline
        \end{tabular}
        
        {\bf Input parameters:} \\
        \begin{tabular}{| l | l | l | l |}
        \hline
        param  & mode & type  & \\
        \hline \hline
        in     & V & any  & input          \\ \hline
        len    & I & int  & optional line fixed delay \\ \hline
        init   & I & any  & optional initialisation \\ \hline
        step   & V & log  & optional line step \\ \hline
        varlen & V & uint & optional line variable delay \\ \hline
        reset  & V & log  & optional line reset \\ \hline
        \end{tabular}
        
        {\bf Parameter=argument allowed:} \\
        Input arguments only.
        
        {\bf Description:} \\
        This module is a delay line implemented using CLB shift registers or
        flip-flops. On each clock cycle the input is written to the line and the
        contents move one item right. The rightmost item appears as the out
        value. This module can function as a sequence generator if the output
        is connected back to the input.
        
        This module is used extensively internally by 3PL to generate single
        cycle delays in execution threads in which context it is implemented as
        a single flip-flop.
        
        For large delays this uses a lot of CLBs and the alternative library
        function {\bf delay()} \autorefph{sec:lfdelay} in library stdlib.3pl
        might be preferable. That library function will call {\bf del()}
        in any case if the delay value is small enough.

        This has 1 to 6 input arguments and 1 output argument. Input arguments
        {\bf len} and {\bf init} are immediate mode. {\bf step} may be
        immediate or target mode as described below. The other input arguments
        must not be modes {\bf immediate}, {\bf queue}, {\bf memory} or {\bf
        clock} and value mode arguments must not contain queue reads.
        
        None of the arguments to this module can contain queue reads. This
        limitation could be easily removed but it is felt that the resulting
        behaviour of this module could become unduly complex. In limited
        circumtances this limitation can be worked around and there are examples
        below to illustrate this.

        The 1st input argument, {\bf in}, is the delay line data input. It can
        be any target type.

        The optional input argument, {\bf len}, specifies a fixed delay.
        This must be an immediate int. If this argument is 0 or omitted and the
        5th argument {\bf varlen} is also omitted, the output will simply be
        connected to the input to give zero delay.

        The optional input argument, {\bf init}, is an initialisation array.
        This  must be immediate mode and an array of the same type as {\bf in}.
        It may be smaller than the delay line length in which case the input end
        of the delay line will be zero initialised. The low order entries in
        the initialisation array are loaded in the output end of the delay line.
        Thus for example a delay of 5 initialised to \{1, 2, 3\} will output
        1 2 3 0 0.
        If initialisation values are too large their high-order bits will be
        truncated.

        The optional input argument, {\bf step}, is a line step enable.
        This  must be type log. If this argument is used the line will only
        advance when {\bf step} is true, otherwise the line will advance on
        every clock cycle.

        The optional input argument, {\bf varlen}, is a variable line
        length. This must be type uint. This argument may be limited or
        unavailable depending on the type of FPGA. For Xilinx the variable
        length is limited to values between 0 and 15 giving delays of 1 to 16
        respectively. A wider argument will be truncated and a narrower one
        zero padded. Note that when the line length is changed data may be
        skipped or repeated. If the line length is changed while stepping is
        inhibited, the line length argument is an address with which the
        contents of the line can be non-destructively examined.

        The optiona input argument, {\bf reset}, returns the
        line to its initial state when true.

        The output argument, {\bf out}, is the delay line data output. This
        argument must be value mode and must be the same type as the data
        input {\bf in}.
        
        If the delay line is initialised, these initial values appear at the
        output during the initial {\em n} line steps, where {\em n} is the
        fixed delay (or the initial delay selected by {\bf varlen}
        if used).
        
        For Xilinx the following apply -
        \blist
        \item   For FPGAs prior to Virtex5 the variable length argument is limited to values
                between 0 and 15 inclusive, giving line lengths of 1 to
                16 respectively. From Virtex5 onwards the lengths are 1 to 32.
        \item   If the reset is used, the variable length argument cannot be
                used.
        \item   If the reset is used or the variable length parameter is not
                used and the delay is 3 or less, the delay line will be
                implemented using flip-flops, otherwise it will be implemented
                using CLB LUT shift registers.
        \blend
        Implementation using CLB LUT shift registers minimises resources,
        requiring only one CLB LUT for every 16-bit delay line.
        
        {\bf Note: the fixed delay parameter {\em len} gives the number of cycles
        to be delayed but the variable delay parameter {\em varlen} is the
        number of cycles delay minus one!}      
            
        In Xilinx devices there is a substantial delay between the clock edge
        and the arrival of the shift register signal at the output of the logic
        block (2.8nS to 3.5nS for a -6 speed XC2VP and about 1.15 to 1.43nS
        for  a -2 speed XC5V). The delay for a flip-flop is very much smaller.
        For this reason a delay line with fixed length is implemented using a
        flip-flop for the last delay stage. If a variable length delay line is
        required, this flip-flop is not inserted and hence the significant
        output delay must be considered. It is preferable in that case to
        assign the delay line output directly to a variable, i.e. with no
        intervening combinatorial logic, to avoid introducing any extra path
        delay.
        
        The limitation stated above that inputs cannot contain queue reads
        can be overcome in some cases. For example to use a queue as data input
        we can use the examine operator as long as a queue wait controls the
        delay step. For example to form a running average of length 8 -

        \begin{lstlisting}[gobble=12, language=threepl]
            queue(q, "uint:8");
            static(rsum, "uint:11");
            queue(rav, "uint:8");
            value(out, "uint:8");
            
            del(out <- ?q, 8,, <q);     // there is no queue read here!
            
            seq {
                rsum <- rsum + q - out; // there is a queue read here
                rav <<- rsum / 8;
            }
        \end{lstlisting}
        Here a data stream is written to queue {\bf q} and a running average
        appears as a stream through queue {\bf rav}. The data input does not
        contain a queue read as the examine operator simply looks at the queue
        output, valid or not, and does not pop a datum. It is important to note
        that the read availability of {\bf q} is tested in two places, the
        reading of {\bf q} and the expression {\bf \verb'<q'}. Although {\bf
        del()} is a distinct module {\bf \verb'<q'} is evaluated as an argument
        to {\bf del()} in the current scope, not the scope of {\bf del()}, so
        there is only a single destination to {\bf q}, or rather both
        availability tests are in the same destination module.
        
        Another approach is used in library function {\bf del()} in
        xilinx/common.3pl \autorefph{sec:lfdel} which calls inbuilt module
        {\bf del()}.
        


       


    \section{resync - resynchronise a level}\label{sec:imresync}
        
        {\bf Call:} \\
        resync(out \verb'<-' in);
        
        {\bf Output parameters:} \\
        \begin{tabular}{| l | l | l | l |}
        \hline
        param & mode & type & \\
        \hline \hline
        out & V & log & output pulse \\ \hline
        \end{tabular}
        
        {\bf Input parameters:} \\
        \begin{tabular}{| l | l | l | l |}
        \hline
        param & mode & type & \\
        \hline \hline
        in    & V SV S PR & log & target value pulse or level \\ \hline
        \end{tabular}
        
        {\bf Parameter=argument allowed:} \\
        No
        
        {\bf Description:} \\
        {\bf in} is an expression or a value, selectvalue, static, input or
        priority variable of type log.
        
        {\bf out} is a value of type log.

        When the input rises and is sampled by the clock associated
        with that variable, the output will produce a single cycle pulse
        synchronised with the current clock domain.

        If the input is an input, selectvalue, static or priority mode variable
        and it has not yet been referenced or assigned then it will be
        given the current clock domain. An unassigned value variable
        will also be given the current clock domain. A value variable that
        has been assigned but has a null clock domain will result in error
        termination.

